%% Results (MedIA revision, v2 pipeline).
%% Every number is read from results/v2/analysis/{C,S}/... (snapshot 1381bdd759f9; ID-validation
%% selection unless stated), from the tables generated by src/v2/paper_tables.py (paper/media/tables/)
%% or from results/v2/analysis/paper_numbers.json, which the same script writes.
\section{Results}
\label{sec:results}

No model is trained on labeled data from its test domain. Apart from the
test-time adaptation arms, which use unlabeled test images, every result below
is the performance of a model applied unchanged to an unseen site.
Unless stated otherwise, results use the primary model-selection rule
(training-domain validation) and test AUROC. We lead with the standard tier,
whose models are closer to current practice, and give the compact tier
alongside. Arms marked $^{\dagger}$ use unlabeled test-domain images and are
not like-for-like competitors of the other arms
(Section~\ref{sec:setup-arms}).

\subsection{Performance at every held-out domain}
\label{sec:results-domains}

Figure~\ref{fig:perdomain} shows the seed-mean test AUROC of every
standard-tier arm at every held-out domain, and Table~\ref{tab:main-S}
summarizes it by the worst domain and the mean. Averages hide large
differences between domains. Standard-tier ERM has a mean AUROC of 0.785 on
the canine cohort but scores 0.534 on the 3DHistech P1000 scanner, which is
close to chance. That scanner is the worst domain for 17 of the 21 arms. On
Camelyon17, ERM averages 0.955 and drops to 0.895 at center~2. The two mitosis
cohorts are flatter: ERM ranges from 0.847 to 0.879 on MIDOG~2021 and from
0.870 to 0.927 on MIDOG++. The compact tier shows the same pattern at a lower
level (Supplementary Table~\ref{S-tab:main-C}). There, ERM's worst domains score
0.698 (Camelyon17), 0.551 (canine), 0.765 (MIDOG~2021) and 0.766 (MIDOG++).

The worst domain is not the same for every arm. On Camelyon17, center~2 is the
worst domain for 10 of 21 standard-tier arms, and the others fail at
centers~1, 3 or 4. A study that reports a single held-out hospital therefore
measures one point of a distribution whose shape depends on the method.

\begin{figure}[pos=tbp]
  \centering
  \includegraphics[width=\linewidth]{fig_perdomain_S}
  \caption{Test AUROC of every standard-tier arm at every held-out domain
  (mean over three seeds; ID-validation selection). Outlined cells mark each
  arm's worst domain. Domains: Camelyon17 hospitals H0--H4; canine scanners;
  MIDOG~2021 scanners; MIDOG++ tumor types ((c) canine, (h) human).
  $^{\dagger}$Transductive: uses the unlabeled images of the test domain at
  prediction time. The compact tier is in Supplementary
  Fig.~\ref{S-fig:perdomain-C}.}
  \label{fig:perdomain}
\end{figure}

\begin{table}[width=\linewidth,pos=p]
\caption{Standard tier: worst-domain AUROC (rank among the 21 arms) and mean
AUROC over held-out domains, ID-validation selection, three seeds. Bold: best
worst-domain AUROC in the cohort. $^{\dagger}$Transductive. FT: fine-tuned;
BT: Barlow Twins; IN: ImageNet. The compact tier is in Supplementary
Table~\ref{S-tab:main-C}.}
\label{tab:main-S}
\footnotesize
\setlength{\tabcolsep}{3pt}
% generated by src/v2/paper_tables.py from results/v2/analysis/S/*/summary_id.csv, ranks_id.csv
\begin{tabular*}{\tblwidth}{@{\extracolsep{\fill}}l*{4}{rr}@{}}
\toprule
Arm & \multicolumn{2}{c}{Camelyon17 ($K=5$)} & \multicolumn{2}{c}{Canine cSCC ($K=5$)} & \multicolumn{2}{c}{MIDOG 2021 ($K=3$)} & \multicolumn{2}{c}{MIDOG++ ($K=7$)} \\
\cmidrule(lr){2-3}\cmidrule(lr){4-5}\cmidrule(lr){6-7}\cmidrule(lr){8-9}
 & worst (rank) & mean & worst (rank) & mean & worst (rank) & mean & worst (rank) & mean \\
\midrule
\multicolumn{9}{@{}l}{\emph{Reference and controls}}\\
ERM & 0.895 (14) & 0.955 & 0.534 (21) & 0.785 & 0.847 (10) & 0.860 & 0.870 (8) & 0.897 \\
\addlinespace[2pt]
\multicolumn{9}{@{}l}{\emph{Domain-aware training}}\\
GroupDRO & 0.893 (16) & 0.953 & 0.606 (19) & 0.803 & 0.826 (16) & 0.848 & 0.838 (16) & 0.890 \\
IRM & 0.819 (19) & 0.923 & 0.559 (20) & 0.776 & 0.837 (13) & 0.850 & 0.638 (20) & 0.861 \\
CORAL & 0.923 (9) & 0.966 & 0.702 (15) & 0.823 & 0.820 (18) & 0.846 & 0.874 (6) & 0.899 \\
MixStyle & 0.792 (20) & 0.924 & 0.632 (16) & 0.812 & 0.846 (11) & 0.864 & 0.840 (15) & 0.891 \\
Fish & 0.838 (18) & 0.911 & 0.792 (11) & 0.844 & 0.820 (17) & 0.845 & 0.775 (17) & 0.871 \\
LISA & 0.886 (17) & 0.948 & 0.734 (14) & 0.828 & 0.847 (9) & 0.861 & 0.862 (14) & 0.888 \\
\addlinespace[2pt]
\multicolumn{9}{@{}l}{\emph{Stain normalization and color}}\\
Macenko$^{\dagger}$ & 0.915 (10) & 0.962 & 0.819 (10) & 0.881 & 0.860 (6) & 0.872 & 0.873 (7) & 0.906 \\
Macenko (per tile) & 0.557 (21) & 0.891 & 0.783 (12) & 0.854 & 0.846 (12) & 0.867 & 0.877 (5) & 0.895 \\
H\&E jitter & 0.958 (4) & 0.979 & 0.850 (7) & 0.903 & 0.875 (3) & 0.888 & 0.880 (3) & 0.903 \\
TIA-style$^{\dagger}$ & 0.934 (6) & 0.970 & 0.823 (9) & 0.887 & 0.871 (4) & 0.881 & 0.877 (4) & 0.905 \\
\addlinespace[2pt]
\multicolumn{9}{@{}l}{\emph{Self-supervision}}\\
SimCLR & 0.929 (7) & 0.968 & 0.607 (18) & 0.816 & 0.836 (14) & 0.855 & 0.882 (2) & 0.901 \\
\addlinespace[2pt]
\multicolumn{9}{@{}l}{\emph{Sharpness-aware}}\\
SAM & 0.900 (13) & 0.960 & 0.624 (17) & 0.818 & 0.851 (8) & 0.872 & \textbf{0.884} (1) & 0.907 \\
\addlinespace[2pt]
\multicolumn{9}{@{}l}{\emph{Test-time adaptation}}\\
BN-adapt (ERM)$^{\dagger}$ & 0.929 (8) & 0.970 & 0.882 (2) & 0.905 & 0.870 (5) & 0.882 & 0.868 (9) & 0.893 \\
\addlinespace[2pt]
\multicolumn{9}{@{}l}{\emph{Pretrained representations}}\\
Lunit BT FT & 0.895 (15) & 0.954 & 0.833 (8) & 0.885 & 0.877 (2) & 0.885 & 0.863 (12) & 0.891 \\
Lunit BT FT + H\&E jitter & 0.902 (12) & 0.957 & 0.866 (5) & 0.913 & \textbf{0.885} (1) & 0.891 & 0.868 (10) & 0.896 \\
Lunit DINO FT & \textbf{0.986} (1) & 0.989 & 0.859 (6) & 0.924 & 0.829 (15) & 0.854 & 0.862 (13) & 0.898 \\
Lunit DINO FT + H\&E jitter & 0.982 (2) & 0.991 & 0.871 (3) & 0.929 & 0.858 (7) & 0.875 & 0.867 (11) & 0.904 \\
RN50-IN probe & 0.906 (11) & 0.948 & 0.772 (13) & 0.820 & 0.623 (21) & 0.647 & 0.571 (21) & 0.614 \\
Lunit BT probe & 0.965 (3) & 0.977 & 0.870 (4) & 0.896 & 0.790 (19) & 0.800 & 0.739 (18) & 0.802 \\
Lunit DINO probe & 0.952 (5) & 0.979 & \textbf{0.890} (1) & 0.931 & 0.736 (20) & 0.746 & 0.712 (19) & 0.756 \\
\bottomrule
\end{tabular*}

\end{table}

\subsection{Where the variance of a held-out score comes from}
\label{sec:results-variance}

Figure~\ref{fig:variance} and Table~\ref{tab:evidence} report the variance
decomposition of Section~\ref{sec:protocol-variance}. For most arms, most of
the variance of a single run is due to which domain is held out. The median
ICC over arms lies between 0.72 and 0.95 in every cohort and at both tiers.
It is highest on the canine cohort and MIDOG++ at the standard tier (0.95 in
both), where 20 of 21 and 18 of 21 arms have an ICC interval whose lower limit
exceeds 0.5. For comparison, the same estimators give a median ICC of 0.98 on
DomainBed's published results outside TerraIncognita
(Supplementary Section~\ref{S-app:domainbed}).

Run-to-run noise is nevertheless far from negligible, and it is largest at
the hard domains. At the compact tier, the five ERM seeds at Camelyon17
center~2 span 0.599 to 0.929 AUROC, and at center~4, ERM's worst domain, they
span 0.592 to 0.797 (Supplementary Fig.~\ref{S-fig:seedspread}). At the standard
tier, the three ERM seeds on the canine P1000 scanner span 0.461 to 0.598. For
some arms the spread at the worst domain is larger still: the seeds of
standard-tier IRM at Camelyon17 center~3 span 0.50 to 0.98, and those of
compact-tier per-tile Macenko at center~4 span 0.41 to 0.74.

The intervals are wide. On MIDOG~2021, with $K=3$, only 3 of 21 standard-tier
arms have an ICC lower limit above 0.5, and for five arms the lower limit is
zero. Estimates at the zero boundary occur for three standard-tier arms on
Camelyon17, and the Bayesian posterior gives them one-sided intervals
(Fig.~\ref{fig:variance}a). The ICC is therefore best read as a coarse
description of each cohort, not as a property of a method.

\begin{figure}[pos=tbp]
  \centering
  \includegraphics[width=0.92\linewidth]{fig_variance_S}
  \caption{Variance decomposition at the standard tier. (a) ICC(1), the share
  of single-run variance due to the held-out domain, with the REML estimate and
  exact-$F$ 95\% interval and the Bayesian (half-Cauchy) posterior median and
  95\% credible interval; one-sided $[0,q_{0.95}]$ at the zero boundary.
  (b) Between-domain SD $\hat\sigma_D$ against between-run SD $\hat\sigma_R$
  per arm; rays mark constant ICC. $^{\dagger}$Transductive. The compact tier
  is in Supplementary Fig.~\ref{S-fig:variance-C}.}
  \label{fig:variance}
\end{figure}

\begin{table}[width=\linewidth,pos=t]
\caption{Summary of the evidence per tier and cohort (ID-validation selection).
ICC: median and interquartile range over arms, and the number of arms whose
exact 95\% ICC interval has a lower limit above 0.5. ERM
$\hat\sigma_D$/$\hat\sigma_R$: between-domain and between-run SD of ERM.
$p<0.05$ and Holm: number of arms whose paired mean difference from ERM is
significant before and after Holm adjustment, split into better ($+$) and worse
($-$). Largest mean gain: the arm with the largest mean AUROC difference from
ERM. $\hat\sigma_\Delta$: median over arms of the SD of the per-domain
difference from ERM with the tier's seed count (Eq.~\ref{eq:diffvar}).
$K_{0.05}$: median over arms of the number of domains needed for 80\% power to
detect a mean gain of 0.05 AUROC, for independent folds and, in brackets, under
fold dependence $\rho=0.5$ (Eq.~\ref{eq:power}). Seeds: 3 (S), 5 (C).
$^{\dagger}$Transductive.}
\label{tab:evidence}
\footnotesize
\setlength{\tabcolsep}{2.5pt}
% generated by src/v2/paper_tables.py from variance_id.csv, contrasts_id.csv, power_id.csv
\begin{tabular*}{\tblwidth}{@{\extracolsep{\fill}}llrlrlrr>{\raggedright\arraybackslash}p{2.2cm}rr@{}}
\toprule
Tier & Cohort & $K$ & ICC median [IQR] & CI$_{\text{lo}}>0.5$ & ERM $\hat\sigma_D$ / $\hat\sigma_R$ & $p<0.05$ ($+$/$-$) & Holm ($+$/$-$) & Largest gain over ERM & $\hat\sigma_\Delta$ & $K_{0.05}$ \\
\midrule
S & C17 & 5 & 0.73 [0.40, 0.93] & 7/21 & 0.036 / 0.018 & 0 / 0 & 0 / 0 & Lunit DINO FT + H\&E jitter (+0.036) & 0.036 & 7 (14) \\
 & Canine & 5 & 0.95 [0.92, 0.98] & 20/21 & 0.149 / 0.034 & 0 / 0 & 0 / 0 & Lunit DINO probe (+0.145) & 0.126 & 53 (137) \\
 & MIDOG21 & 3 & 0.76 [0.71, 0.85] & 3/21 & 0.016 / 0.010 & 4 / 3 & 0 / 0 & Lunit BT FT + H\&E jitter (+0.031) & 0.013 & 3 (4) \\
 & MIDOG++ & 7 & 0.95 [0.90, 0.97] & 18/21 & 0.020 / 0.005 & 2 / 4 & 0 / 3 & SAM (+0.010) & 0.012 & 3 (5) \\
\addlinespace[3pt]
C & C17 & 5 & 0.81 [0.67, 0.89] & 11/20 & 0.101 / 0.080 & 0 / 1 & 0 / 0 & TIA-style$^{\dagger}$ (+0.107) & 0.094 & 30 (76) \\
 & Canine & 5 & 0.80 [0.68, 0.88] & 10/20 & 0.103 / 0.055 & 0 / 0 & 0 / 0 & BN-adapt (ERM)$^{\dagger}$ (+0.123) & 0.103 & 36 (92) \\
 & MIDOG21 & 3 & 0.72 [0.38, 0.83] & 3/20 & 0.037 / 0.025 & 7 / 0 & 0 / 0 & BN-adapt (H\&E jitter)$^{\dagger}$ (+0.091) & 0.028 & 5 (8) \\
 & MIDOG++ & 7 & 0.85 [0.79, 0.92] & 16/20 & 0.033 / 0.014 & 5 / 7 & 2 / 3 & Compute-matched ERM (+0.060) & 0.023 & 4 (10) \\
\bottomrule
\end{tabular*}

\end{table}

\subsection{Can the differences between methods be detected?}
\label{sec:results-contrasts}

Table~\ref{tab:evidence} counts the paired comparisons with ERM. At the
standard tier, no arm is significantly better than ERM after Holm adjustment in
any cohort. The only Holm-significant differences are three frozen linear
probes that are worse than ERM on MIDOG++. At the compact tier, two arms
improve on ERM after adjustment, both on MIDOG++: compute-matched ERM ($+0.060$)
and per-slide Macenko$^{\dagger}$ ($+0.041$). The augmentation-only control,
LISA and IRM are significantly worse there.

Lack of significance does not mean that the differences are small. On the
canine cohort, the Lunit DINO frozen probe improves mean AUROC over ERM by
0.145 and is better at all five scanners, yet its unadjusted $p$ is 0.071.
With five domains, even an exact sign-flip test cannot go below 0.0625 when
every domain improves. On the tumor cohorts the gains are large but vary
strongly from domain to domain. On the mitosis cohorts the reverse holds: the
differences are measured precisely but are small. The largest standard-tier
gain on MIDOG++ is $+0.010$ (SAM, better at all seven domains, Holm
$p=0.053$).

The power analysis puts numbers on this (Fig.~\ref{fig:power},
Table~\ref{tab:evidence}). The SD of the per-domain difference from ERM,
$\hat\sigma_\Delta$, has a median of 0.012--0.028 on the mitosis cohorts but
0.036--0.126 on the tumor cohorts. For a median arm to detect a mean gain of
0.05 AUROC with 80\% power, the standard tier would need 7 Camelyon17
hospitals and 53 canine scanners. The compact tier would need 30 and 36. If
the folds are correlated through their shared training domains
($\rho=0.5$), these numbers roughly double or triple. On the mitosis cohorts,
3--5 domains suffice for a gain of that size, but the observed gains there are
much smaller. Adding seeds helps little once $\sigma_\Delta$ is dominated by
the domain-by-arm interaction (Fig.~\ref{fig:power}). Domains, not seeds, are
the scarce resource.

\begin{figure}[pos=tbp]
  \centering
  \includegraphics[width=\linewidth]{fig_power_S}
  \caption{Number of held-out domains $K$ needed for 80\% power to detect a
  mean AUROC gain $\delta$ over ERM (two-sided paired $t$ test,
  $\alpha=0.05$), at the standard tier, for the training arm with the median
  $\sigma_\Delta$ in each cohort. Solid: independent folds; dashed: fold
  dependence $\rho=0.5$; dotted: $K$ of this study. Colors: seeds per arm. The
  compact tier is in Supplementary Fig.~\ref{S-fig:power-C}.}
  \label{fig:power}
\end{figure}

\subsection{Rankings within a cohort, and the effect of model scale}
\label{sec:results-ranks}

Figure~\ref{fig:ranks}a gives each arm's worst-domain rank with bootstrap
intervals over seeds and over domains. Resampling domains usually widens the
intervals more than resampling seeds. The median width of the domain-bootstrap
interval is 3 to 7 ranks out of 21 at the standard tier, against 2 to 6 for
the seed bootstrap. Only one standard-tier arm has the rank interval $[1,1]$ under both
bootstraps: the Lunit DINO probe on the canine cohort, which ranks first in
99\% of domain resamples. The next most certain is fine-tuned Barlow Twins with
H\&E jitter on MIDOG~2021 (96\%, interval $[1,3]$). On MIDOG++, the top arm
(SAM) ranks first in only 45\% of domain resamples.

A few patterns hold across cohorts.
\begin{itemize}
  \item \emph{Pathology foundation models lead on the tumor cohorts.} The
    fine-tuned Lunit DINO ViT ranks first on Camelyon17 (worst domain 0.986),
    and its frozen probe ranks first on the canine cohort (0.890). On the
    mitosis cohorts, the frozen probes fall to the bottom: the DINO probe ranks
    20th and 19th, and the ImageNet ResNet-50 probe last on both. Fine-tuning
    narrows the gap: fine-tuned Barlow Twins with H\&E jitter ranks first on
    MIDOG~2021 (0.885), and the fine-tuned DINO ViT ranks 15th and 13th on the
    two mitosis cohorts.
  \item \emph{H\&E color jitter is consistently good.} It ranks 4th, 7th, 3rd
    and 3rd at the standard tier and 3rd, 8th, 5th and 8th at the compact
    tier.
  \item \emph{The domain-generalization objectives rarely help.} Under
    ID-validation selection, none of the six (GroupDRO, IRM, CORAL, MixStyle,
    Fish and LISA) ranks in the top five on any cohort at either tier. The
    best of them is CORAL on MIDOG++ (6th at the standard tier). Other rules
    can lift a single objective: with oracle selection, CORAL ranks 4th on
    MIDOG++ at the standard tier and LISA 2nd on the canine cohort at the
    compact tier.
  \item \emph{Transductive arms are strong.} BN-adapt on ERM ranks 2nd on the
    canine cohort, and at the compact tier BN-adapt on H\&E jitter ranks among
    the top six everywhere. Both use unlabeled test-domain images.
\end{itemize}

Model scale changes the picture (Fig.~\ref{fig:scale}). Moving from the
compact CNN to ImageNet ResNet-50 raises ERM's worst-domain AUROC on
Camelyon17 (0.698 to 0.895) and on the mitosis cohorts (0.765 to 0.847;
0.766 to 0.870), but not on the canine cohort (0.551 to 0.534). For the 14
arms run at both tiers, Kendall's $\tau_b$ between the two tiers' worst-domain
rankings is 0.60 (Camelyon17), 0.63 (canine), 0.32 (MIDOG~2021) and 0.30
(MIDOG++). Conclusions drawn from the compact model therefore transfer only in
part. Compute-matched ERM, the best compact arm on MIDOG++, was not run at the
standard tier, so that result cannot be checked there.

Hyperparameters matter as much as methods. The Barlow Twins fine-tunes first
inherited the learning rate tuned for ImageNet ResNet-50, and training failed
on some folds. Tuning their own rate raised mean off-site AUROC on MIDOG++
from 0.715 to 0.891 (Supplementary Section~\ref{S-app:tuning-bt}). Stain normalization shows
the same sensitivity to implementation. Estimating the Macenko stain matrix
per tile instead of per slide drops the arm from rank~10 to last on
Camelyon17 at the standard tier, because it fails at one center
(0.557 against 0.976 for the per-slide estimate;
Supplementary Section~\ref{S-app:macenko}).

\begin{figure}[pos=tbp]
  \centering
  \includegraphics[width=0.85\linewidth]{fig_ranks_S}
  \caption{Rankings at the standard tier. (a) Rank by worst-domain AUROC
  (1 = best) with 95\% intervals from bootstrapping seeds within domains
  (thick) and domains (thin). (b) Ranks of all 21 arms across the four cohorts;
  ERM, SimCLR, H\&E jitter and CORAL are highlighted. $^{\dagger}$Transductive.
  The compact tier is in Supplementary Fig.~\ref{S-fig:ranks-C}.}
  \label{fig:ranks}
\end{figure}

\begin{figure}[pos=tbp]
  \centering
  \includegraphics[width=\linewidth]{fig_scale}
  \caption{Effect of model scale and pre-training. Worst-domain AUROC (filled,
  with seed-bootstrap 95\% interval) and mean AUROC (open) for ERM, the best
  domain-generalization objective, H\&E jitter and SimCLR at the compact tier
  (compact CNN, 64~px) and the standard tier (ImageNet ResNet-50, 96~px), and
  for the pathology foundation models. The $x$-axes are scaled per cohort.}
  \label{fig:scale}
\end{figure}

\subsection{Do rankings follow the task or the type of shift?}
\label{sec:results-cross}

The arm-by-cohort interaction is significant at both tiers (wild bootstrap,
$p<0.001$ at the standard tier and $p=0.004$ at the compact tier): methods do
not rank the same way on every cohort. Figure~\ref{fig:kendall} shows how the
rankings agree between pairs of cohorts.

At the standard tier, agreement follows the task. The two tumor cohorts agree
with $\tau_b=0.43$ (95\% interval [0.28, 0.52]) and the two mitosis cohorts
with $\tau_b=0.42$ [0.30, 0.50]. The pairs that share only the shift type agree
much less (0.15 and 0.14), and the pairs that share neither do not agree at all
(0.03 and $-0.06$). The mean $\tau_b$ of same-task pairs exceeds that of
same-shift pairs by 0.28 [0.19, 0.40]. In the $2\times2$ decomposition of
Eq.~\eqref{eq:2x2}, the task contrast accounts for 84\% of the between-cohort
sum of squares, the shift contrast for 10\% and their interaction for 7\%. The
pattern is the same under the last-checkpoint (85\% task) and oracle (90\%)
rules.

At the compact tier the picture is less clear. The task contrast accounts for
51\% and the shift contrast for 32\%, and same-task pairs agree only slightly
more than same-shift pairs (difference 0.06 [$-0.04$, 0.14]). The canine and
MIDOG~2021 cohorts, which share the shift type but not the task, agree with
$\tau_b=0.54$, almost as much as the two mitosis cohorts (0.59).

Two cautions apply. First, with four cohorts there are only three ways to
split them into two pairs, and the task split ranks first among them at both
tiers. A permutation test at the cohort level cannot go below $p=1/3$. The
decomposition describes the interaction; it cannot prove that rankings are
organized by task. Second, the 150 breast-cancer images of MIDOG++ are the
MIDOG~2021 images, so the two mitosis cohorts share one of MIDOG++'s seven
domains. This overlap may raise their agreement. The tumor cohorts share no
images, and their agreement at the standard tier is as high.

\begin{figure}[pos=tbp]
  \centering
  \includegraphics[width=\linewidth]{fig_kendall_S}\\[4pt]
  \includegraphics[width=\linewidth]{fig_kendall_C}
  \caption{Agreement of rankings between cohorts at the standard tier (top) and
  the compact tier (bottom). (a) Kendall's $\tau_b$ of worst-domain rankings
  with 95\% seed-bootstrap intervals. (b) Split of the between-cohort sum of
  squares of the per-domain gains over ERM into task, shift type and their
  interaction (Eq.~\ref{eq:2x2}); descriptive, since only three two-versus-two
  splits exist. (c) $\tau_b$ by type of cohort pair.}
  \label{fig:kendall}
\end{figure}

\subsection{Sensitivity analyses}
\label{sec:results-sensitivity}

\paragraph{Model-selection rule} Table~\ref{tab:selection} compares the
worst-domain ranking under each rule with the ID-validation ranking. At the
standard tier, the rules agree closely ($\tau_b$ 0.83--0.99), and the top arm
is the same under every rule on Camelyon17 and MIDOG++. At the compact tier,
agreement is weaker where run noise is large. On the canine cohort,
$\tau_b$ is 0.38 for OOD-validation selection and 0.32 for the oracle, and on
Camelyon17 it is 0.53--0.63. A compact-tier ranking thus depends on the
selection rule as well as on the seeds (Supplementary Fig.~\ref{S-fig:selection-C}).

\begin{table}[width=\linewidth,pos=t]
\caption{Agreement between model-selection rules: Kendall's $\tau_b$ between
the worst-domain ranking under each rule and the ranking under ID-validation
selection, over the arms analyzed under every rule (TTA pseudo-arms are
evaluated only under ID-validation selection). MIDOG~2021 has no OOD-validation
domain. Same top arm: whether every rule puts the same arm first.}
\label{tab:selection}
\footnotesize
% generated by src/v2/paper_tables.py from summary_{id,ood,last,oracle}.csv (worst-domain AUROC)
\begin{tabular*}{\tblwidth}{@{\extracolsep{\fill}}llrrrrl@{}}
\toprule
Tier & Cohort & Arms & $\tau_b$ OOD-val & $\tau_b$ last & $\tau_b$ oracle & Same top arm \\
\midrule
S & Camelyon17 & 20 & 0.84 & 0.83 & 0.85 & yes \\
 & Canine cSCC & 20 & 0.89 & 0.97 & 0.87 & no \\
 & MIDOG 2021 & 20 & -- & 0.85 & 0.85 & no \\
 & MIDOG++ & 20 & 0.97 & 0.99 & 0.96 & yes \\
\addlinespace[3pt]
C & Camelyon17 & 16 & 0.58 & 0.63 & 0.53 & no \\
 & Canine cSCC & 16 & 0.38 & 0.90 & 0.32 & yes \\
 & MIDOG 2021 & 16 & -- & 0.92 & 0.85 & yes \\
 & MIDOG++ & 16 & 1.00 & 1.00 & 0.98 & yes \\
\bottomrule
\end{tabular*}

\end{table}

\paragraph{Scale normalization of MIDOG~2021} Without normalizing the
$\approx$11\% difference in scanner magnification, the compact-tier ranking of
20 arms changes little ($\tau_b=0.78$ for the worst domain and 0.83 for the
mean). No arm's worst-domain AUROC moves by more than 0.039 (Supplementary
Fig.~\ref{S-fig:sensitivity}).

\paragraph{PatchCamelyon} Camelyon17 models were also scored on the
PatchCamelyon test split. PatchCamelyon comes from Camelyon16, whose two
hospitals also contribute to Camelyon17
\citep{veeling2018pcam,litjens2018camelyon}, so it is a secondary check, not an
independent cohort. At the standard tier the Lunit models,
fine-tuned or frozen, score highest, led by fine-tuned DINO with H\&E jitter
(mean over folds 0.908, against 0.772 for ERM). At the compact tier TIA-style
(0.802) and the augmentation arms score highest, against 0.698 for ERM
(Supplementary Table~\ref{S-tab:pcam}).

\paragraph{Calibration} On every cohort at the standard tier, one of the two
fine-tuned foundation models with H\&E jitter has the lowest mean ECE
(0.027--0.095), against
0.089--0.180 for ERM (Supplementary Fig.~\ref{S-fig:calibration-S}). Per-slide
normalization with jitter (TIA-style) and jitter alone are close: on the
canine cohort, their ECE is 0.075 and 0.080 at the standard tier and 0.047 and
0.065 at the compact tier.

\paragraph{Frozen representations} Among nine frozen backbones under one
linear-probe protocol (Supplementary Fig.~\ref{S-fig:ladder}), the Lunit Barlow Twins
ResNet-50 has the best worst-domain AUROC averaged over cohorts (0.841),
followed by the Lunit DINO ViT-S/16 (0.822). Pathology pre-training is not
enough by itself, however. The Lunit SwAV (0.707) and MoCo~v2 (0.734)
ResNet-50s are below ImageNet ConvNeXt-T (0.780) and ViT-B/16 (0.774), and
SwAV is below even the ImageNet ResNet-50 (0.718).

\subsection{The site acceptance test}
\label{sec:results-sat}

Figure~\ref{fig:sat} and Table~\ref{tab:sat} report the validation of the
site acceptance test of Section~\ref{sec:sat} over the 1,200 standard-tier
runs, with every held-out site playing the new site in turn. The development
prior alone is not a safe basis for a decision. Without local labels it
accepted 8.6\% of the sites whose true AUROC was below 0.80, because a model
that works at the development sites can still fail at a new one.

With local labels, the test kept false acceptance at or below 2.1\% for every
target and sample size, and at or below 1.3\% for sites more than 0.02 AUROC
below the target. Its 90\% credible intervals covered the true site AUROC in
92--93\% of cases. For targets of 0.80 and 0.85 it needed about half as many labels as local
labels alone. At a target of 0.80, it accepted 75\% of the good sites with 50 labels and 84\% with
100, while local labels alone accepted 48\% and 70\% and needed 200 labels to
reach 84\%. At a target of 0.85, the test reached with 100 labels (63\%) what
local labels alone reached with 200. Its point estimates were also more
accurate, with a mean absolute error of 0.024 AUROC at 50 labels against 0.036.
The price is a slightly higher false-acceptance rate than local labels alone,
by about one percentage point at the lower targets. The gain is largest where
the development sites score well above the target, as on Camelyon17 and
MIDOG++. Where they score near or below it, as on the canine cohort and at the
compact tier, the prior is pessimistic and can make the test less powerful
than local labels alone (Supplementary Table~\ref{S-tab:sat-cohorts}).

The advantage has limits. At the strict target of 0.90 with 200 or more
labels, local labels alone were slightly more permissive (55\% against 52\%
accepted at 200 labels). In the hardest subgroup, Camelyon17 at a target of
0.90, where most runs score above 0.95 and the 24 runs below the target mostly
fall just short of it, false acceptance reached 8.7\% with 25 labels and 5.7\%
with 50, and fell to 3.3\% at 100 and 1.7\% at 200. An $\epsilon$-contaminated
prior, which gives 10\% of the prior mass to a vague component centered at an
AUROC of 0.5, lowered these values to 5.4\%, 4.5\% and 2.8\%, at a cost in power
at small samples (for example from 38\% to 28\% at 25 labels and a target of
0.85; Supplementary Section~\ref{S-sup:sat}). We added this variant after
seeing the subgroup result and report it as a sensitivity analysis. At the
compact tier, whose sites score lower and differ more, the test was as safe
(false acceptance at most 0.4\%), but it gained over local labels alone only
at small samples and lost some power at larger ones. The decision threshold $1-\alpha$ is a posterior probability, so these
false-acceptance rates are empirical averages over the mix of bad sites in the
study, not a guaranteed error rate; for runs just below the target (within
0.02) false acceptance was higher, up to 5.8\% (target 0.80, 50 labels). The
label planner was close to the observed acceptance at the compact tier and
slightly optimistic at the standard tier, where it predicted 75\% acceptance of good
sites at 100 labels against 68\% observed.

\begin{figure}[pos=tbp]
  \centering
  \includegraphics[width=\linewidth]{fig_sat_S}
  \caption{Validation of the site acceptance test at the standard tier.
  Top: share of good sites (true AUROC $\ge$ target) that each rule accepts;
  bottom: share of bad sites (true AUROC $<$ target) that it accepts, with the
  5\% level dotted. Means over runs of the acceptance rate over 200 random
  labeled samples per run and sample size. The compact tier is in
  Supplementary Fig.~\ref{S-fig:sat-C}.}
  \label{fig:sat}
\end{figure}

\begin{table}[width=\linewidth,pos=tbp]
\caption{Site acceptance test at the standard tier (1,200 runs). FA: false
acceptance, the percentage of runs with true site AUROC below the target $T$
that are accepted; power: the percentage of runs at or above $T$ that are
accepted; $m$: local labeled patches; MAE: mean absolute error of the point
estimate of the site AUROC; cover.: coverage of the 90\% interval. Rule: accept
if $P(\text{AUROC}\ge T)\ge0.95$ (local labels only: one-sided 95\% Wald bound
$\ge T$). The compact tier and the robust variant are in Supplementary
Section~\ref{S-sup:sat}.}
\label{tab:sat}
\footnotesize
\setlength{\tabcolsep}{3pt}
% generated by src/v2/site_acceptance.py tables (results/v2/site_acceptance/sat_rows_*.csv)
\begin{tabular*}{\tblwidth}{@{\extracolsep{\fill}}lrrrrrrrrr@{}}
\toprule
 & & \multicolumn{2}{c}{$T=0.80$} & \multicolumn{2}{c}{$T=0.85$} & \multicolumn{2}{c}{$T=0.90$} & & \\
\cmidrule(lr){3-4}\cmidrule(lr){5-6}\cmidrule(lr){7-8}
Rule & $m$ & FA (\%) & power (\%) & FA (\%) & power (\%) & FA (\%) & power (\%) & MAE & cover. (\%) \\
\midrule
Development prior only & 0 & 8.6 & 53 & 2.8 & 21 & 0.8 & 15 & -- & -- \\
\addlinespace[2pt]
Site acceptance test & 25 & 2.1 & 65 & 0.8 & 38 & 0.3 & 19 & 0.029 & 92 \\
 & 50 & 2.0 & 75 & 0.9 & 51 & 0.3 & 31 & 0.024 & 92 \\
 & 100 & 1.4 & 84 & 0.8 & 63 & 0.2 & 44 & 0.019 & 92 \\
 & 200 & 0.9 & 90 & 0.6 & 72 & 0.2 & 52 & 0.015 & 92 \\
 & 400 & 0.4 & 94 & 0.3 & 79 & 0.1 & 60 & 0.011 & 93 \\
\addlinespace[2pt]
Local labels only & 25 & 0.7 & 11 & 0.0 & 0 & 0.0 & 0 & 0.051 & 94 \\
 & 50 & 1.0 & 48 & 0.6 & 29 & 0.2 & 9 & 0.036 & 94 \\
 & 100 & 0.8 & 70 & 0.6 & 50 & 0.3 & 43 & 0.025 & 93 \\
 & 200 & 0.5 & 84 & 0.4 & 63 & 0.3 & 55 & 0.017 & 93 \\
 & 400 & 0.4 & 92 & 0.2 & 74 & 0.2 & 65 & 0.012 & 94 \\
\bottomrule
\end{tabular*}

\end{table}

% Flush the Results floats before the Discussion.
\clearpage
